\section{Scores, comparison measures, and an amplitude signal}
\label{sec:scores}

A \emph{score} is a real function on the cube of assignments to the bit
positions of the test system. All scores below depend on a bounded number
of positions. Identical functions are identified, even if their sampling
presentations differ. A lookup $h$ is a measurable function from scores
to $\{0,1\}$; measurability means Borel measurability on each finite table
of score values. A raw comparison form is a finite positive measure on
pairs of scores, evaluated as
\[
 \cD(h)=\int |h(B)-h(B')|\dd\mathfrak m(B,B').
\]
Its truth cost at an assignment $x$ is
$\int|B(x)-B'(x)|\dd\mathfrak m(B,B')$.

\begin{theorem}[Analytic score gap]\label{thm:score-gap}
For each fixed $T_0\ge1$, the primitive system of
Section~\ref{sec:events} admits a demand law $\mu$ of scores and a finite
sum $\cD$ of raw comparison forms with the following properties.
All arities, score bounds, numerical dimensions, and total comparison
weight are bounded by constants independent of the input size. The laws
and weights do not depend on a lookup. At a perfect witness $x$,
\[
 \int|B(x)-B'(x)|\dd\mathfrak m(B,B')\le4,
 \qquad
 \mu\{B:B(x)\ge1/4\},\ \mu\{B:B(x)\le-1/4\}\ge10^{-4}.
\]
In a NO instance, every binary measurable lookup with
$\nu=\Var_\mu(h)>0$ satisfies $\cD(h)>T_0\nu$.
Conditional on the discrete sampling data, the numerical laws and score
entries belong to fixed finite families of polynomials with rational
coefficients on rational boxes, with atoms represented as separate cases.
\end{theorem}

We first construct the demand law and comparison measures, with bounded
truth cost at a satisfying assignment. For a putative low-cost coloring,
the resampling comparison supplies an amplitude signal. The remaining
comparisons preserve enough of this signal to obtain the common
conjunction response required by Lemma~\ref{lem:rare-event}.
Section~\ref{sec:paths} constructs that response and completes the proof
of Theorem~\ref{thm:score-gap}.

It suffices to prove the theorem for rational $T_0\ge1$, since one can
replace a given real $T_0$ by a larger rational number. Choose positive
rational mixture weights $\omega_s$ for all primitive types, and write
$P_*$ for their mixture and $p_*=\sum_s\omega_sp_s\in(0,1)$.
A length-$t$ conjunction from the independent mixture law has honest
probability $p=p_*^t$. A fixed pattern of primitive types has honest
probability equal to the product of their $p_s$'s. The ordered conjunction
couplings change one coordinate by one of the prescribed primitive
couplings and use independent, shared, fixed-type samples in all other
coordinates. They give actual pointwise inclusions $r\le r'$.

\subsection{Demand and random half-freezing}

Let $M$ be a positive integer. Stage $j$ contains $n_j$ slots, where
$n_j$ is even. Choose a positive rational coefficient satisfying
\begin{equation}\label{eq:score-scales}
 \frac1{\sqrt{pMn_j}}\le s_j\le\frac2{\sqrt{pMn_j}},\qquad s_j\le1.
\end{equation}
All stage sizes will be chosen below. In each slot independently sample
a primary presentation $\cR_i$ of an iid conjunction $r_i$ and an
independent amplitude with law
\[
 \lambda=\tfrac12\delta_0+\tfrac12\operatorname{Unif}[-1,1].
\]
The demand score and its law are
\[
 B^0=\sum_{j=1}^M s_j\sum_{i\text{ in stage }j}a_i r_i,
 \qquad \mu=\operatorname{Law}(B^0).
\]
At a perfect witness $x$, each summand is symmetric, centered, and bounded
by one, with variance $s_j^2p/6$. Consequently
\[
 \frac16\le\Var(B^0(x))\le\frac23,\qquad
 \E B^0(x)^4
 \le\sum_{j,i}\E(s_ja_ir_i(x))^2+3\Var(B^0(x))^2<4.
\]
For $X=B^0(x)$, Cauchy--Schwarz gives
\[
 \Prob(|X|\ge1/4)
 \ge\frac{(\E X^2-1/16)^2}{\E X^4}
 \ge\frac{25}{9216}.
\]
Symmetry splits this probability equally between the two required tails.

Set $\rho=1/(2048T_0^2)$. Include the form with weight $16T_0$ that
compares $B^0$ with $B^{0,\rho}$, obtained by independently resampling
each complete primary slot with probability $\rho$. Its truth cost is
at most
\[
 16T_0\sqrt{2\rho\Var(B^0(x))}
 \le16T_0\sqrt{4\rho/3}=1/\sqrt6<1.
\]
Throughout the soundness proof suppose, for a contradiction, that
\begin{equation}\label{eq:assumed-score-cut}
 0<\nu=\Var(h(B^0)),\qquad \cD(h)\le T_0\nu.
\end{equation}
In particular,
\begin{equation}\label{eq:resampling-bound}
 \E|h(B^0)-h(B^{0,\rho})|\le\nu/16.
\end{equation}

Independently choose a uniform half of the slots in each stage to freeze.
Let $F_j$ be that half, $I_j$ its complement, and $F=\bigcup_jF_j$.
The information $\cZ$ consists of the freeze sets and all complete
primary variables in frozen slots. Let $\cB$ add to $\cZ$ the primary
presentations of all unfrozen slots, but none of their amplitudes. Define,
for $i\in I_j$,
\[
 s_jb_i^0=\E[h(B^0)\mid\cB,a_i]-\E[h(B^0)\mid\cB],
 \qquad g_i^0=\E[b_i^0\mid\cZ,\cR_i,a_i].
\]
Indices in such formulas are fixed after the freeze sets are specified.
For any family $b=(b_i)$ put
\[
 \|b\|_\Sigma^2=\E\sum_j s_j^2\sum_{i\in I_j}|b_i|^2.
\]

\begin{lemma}[Amplitude signal]\label{lem:amplitude-signal}
Choose $k=\lceil\rho^{-1}\rceil$ and require $n_j\ge4k$. With
$c=4^{-(k+2)}$, Equation~\ref{eq:resampling-bound} implies
\begin{equation}\label{eq:amplitude-signal}
 \|b^0\|_\Sigma^2\le\nu,\qquad
 \sum_j s_j^2\sum_{i\in I_j}\E[|b_i^0|^2\mid\cB]\le\tfrac14,
 \qquad \|g^0\|_\Sigma^2\ge c\nu.
\end{equation}
The constants $k,c$ depend only on $T_0$.
\end{lemma}

\begin{proof}
First consider a single slot. Its erasure observation $O$ is a distinguished
symbol $*$ when $a=0$, and $(\cR,a)$ otherwise. For a square-integrable
function $\psi(O)$, write $c_*=\psi(*)$, let $\psi_1(\cR,u)$ be its
nonerased value, and put $m=\E\psi_1$. Then
\[
 \Var(\E[\psi(O)\mid\cR])
 =\tfrac14\Var(\E_u\psi_1)\le\tfrac14\Var(\psi_1),
 \quad
 \Var(\psi(O))=\tfrac12\Var(\psi_1)+\tfrac14(c_*-m)^2.
\]
Thus the conditional expectation from centered functions of $O$ to
functions of $\cR$ has norm at most $2^{-1/2}$.

The centered function $f=h(B^0)-\E h(B^0)$ is a function of the independent
erasure observations: when a slot has amplitude zero its primary
presentation contributes nothing to the score. Decompose $f$ by active
sets of observation coordinates, using the classical orthogonal
product-space decomposition of Hoeffding and Efron--Stein
\cite{Hoeffding1948,EfronStein1981}. If $P_i$ averages observation $i$, the
component on $S\ne\varnothing$ is
$f_S=\prod_{i\in S}(I-P_i)\prod_{i\notin S}P_if$.
These components are orthogonal and separately centered in their active
coordinates. Conditional expectation onto all primary presentations
preserves this separate centering, so the images of different active
sets remain orthogonal. Applying the one-slot bound in each active
coordinate gives
\begin{equation}\label{eq:erasure-contraction}
 \|\E[f\mid(\cR_i)_i]\|_2^2
 \le\sum_{S\ne\varnothing}2^{-|S|}\|f_S\|_2^2\le\nu/2.
\end{equation}

Refine the decomposition on complete primary slots. Let $\mathsf P_i$
average the whole slot and $\mathsf R_i$ average its amplitude alone.
The operators
\[
 \mathsf C_i=\mathsf P_i,\qquad
 \mathsf D_i=\mathsf R_i-\mathsf P_i,\qquad
 \mathsf A_i=I-\mathsf R_i
\]
are orthogonal projections onto, respectively, constants, centered
functions of primary data, and functions with conditional amplitude mean
zero. Operators in distinct slots commute. Write $f_\eta=\prod_i\eta_if$
for $\eta_i\in\{\mathsf C_i,\mathsf D_i,\mathsf A_i\}$, and let
$S(\eta)=\{i:\eta_i\ne\mathsf C_i\}$ and
$A(\eta)=\{i:\eta_i=\mathsf A_i\}$. Equation~\ref{eq:erasure-contraction}
says that the total squared mass with $A(\eta)\ne\varnothing$ is at
least $\nu/2$. Slot resampling acts on $f_\eta$ by
$(1-\rho)^{|S(\eta)|}$; hence
\[
 \E(h(B^0)-h(B^{0,\rho}))^2
 =2\sum_\eta[1-(1-\rho)^{|S(\eta)|}]\|f_\eta\|_2^2.
\]
Binary differences have equal absolute and squared values. For
$|S(\eta)|>k$, the bracketed coefficient multiplied by two is at least
one. Equation~\ref{eq:resampling-bound} bounds the mass of these
components by $\nu/16$. In particular,
\begin{equation}\label{eq:low-degree-amplitude}
 \sum_{A(\eta)\ne\varnothing,\ |S(\eta)|\le k}
       \|f_\eta\|_2^2\ge3\nu/8.
\end{equation}

Given $\cB$, the unfrozen amplitudes are independent. The variables
$s_jb_i^0$ are the orthogonal projections of
$h(B^0)-\E[h(B^0)\mid\cB]$ onto their separate centered amplitude
spaces. Bessel's inequality bounds their conditional squared norms by
$\Var(h(B^0)\mid\cB)\le1/4$. Averaging and using total variance proves
the first two assertions.

The tower property and amplitude independence give exactly
\[
 s_jg_i^0=\E[h(B^0)\mid\cZ,\cR_i,a_i]
             -\E[h(B^0)\mid\cZ,\cR_i].
\]
For a fixed freeze set, this is the projection
$\mathsf A_i\prod_{\ell\notin F\cup\{i\}}\mathsf P_\ell f$.
After averaging the observed frozen values as well, orthogonality yields
\[
 \|g^0\|_\Sigma^2
 =\sum_\eta\|f_\eta\|_2^2
   \sum_{i\in A(\eta)}
   \Prob_F(i\notin F,\ S(\eta)\setminus\{i\}\subseteq F).
\]
For one component in Equation~\ref{eq:low-degree-amplitude}, distinguish
an $i\in A(\eta)$ in stage $j$, put $d=|S(\eta)|\le k$, and let
$q_\ell$ count its active slots in stage $\ell$. With falling factorials,
the displayed capture probability is
\[
 \frac12\frac{(n_j/2)_{q_j-1}}{(n_j-1)_{q_j-1}}
       \prod_{\ell\ne j}\frac{(n_\ell/2)_{q_\ell}}{(n_\ell)_{q_\ell}}.
\]
Each of its $d-1$ ratio factors is at least $1/4$: its numerator is
at least $n_\ell/2-k+1\ge n_\ell/4$ and denominator at most $n_\ell$.
The probability is therefore at least $\tfrac12 4^{-(d-1)}\ge4^{-k}$.
Keeping just the distinguished contribution and using
Equation~\ref{eq:low-degree-amplitude} gives the claimed $c\nu$.
\end{proof}

We must turn this amplitude signal into the response on actual conjunction
predicates required by Lemma~\ref{lem:rare-event}. We will first construct
a real response, bound its size under each type pattern, and control its
decrease when $r$ is replaced by a containing predicate $r'$. Three kinds
of comparison path provide these estimates. Paths between opposite
contributions of one conjunction bound the response under each pattern;
paths using $r\le r'$ control its order; paths through the random-amplitude
contribution carry the signal to that response. The comparison measures
sample their path orientations independently of the lookup. A common
orientation extracted from the lookup will enter only the analysis.
The estimates must retain errors proportional to $\nu$ throughout, so that
thresholding the response can meet the rare-event criterion even when the
cut variance is arbitrarily small.

\subsection{The order of the constants}

We first fix the signal-loss tolerance $e$, the clipping bound $B_*$, and
the signal threshold $\gamma$, together with the resulting inputs
$b_0,K_0,B_0$ to Lemma~\ref{lem:rare-event}. Choose rational constants
\begin{equation}\label{eq:preconjunction-constants}
 \begin{gathered}
 c_1=c/4,\quad 0<e\le c/160000,\quad
 B_*\ge\max(1,3200(T_0+1)/c),\quad K_*=32(T_0+2),\\
 0<\gamma\le c_1/(16K_*),\quad
 b_*=c_1/(128B_*^2),\quad
 b_0=b_*/[8(T_0+1)],\quad K_0=16/\gamma,\quad B_0=B_*.
 \end{gathered}
\end{equation}
These constants have no dependence on conjunction length. Apply
Lemma~\ref{lem:rare-event} to $b_0,K_0,B_0$ and choose its length $t$
and tolerance $\eps>0$. We can decrease $\eps$ to make it rational.
From now on $p=p_*^t$ is fixed. Let $V$ be the number of fixed type
patterns together with the iid conjunction law, and let $Q$ be the
number of required ordered conjunction couplings, counting the changed
coordinate. Our path menu below has $L=V+2Q+2$ entries. Set
$A_t=16V+80Q$, and choose rational $\Delta>0$ so small that
\begin{equation}\label{eq:delta-choice}
 \Delta\le b_*p\gamma/160,\qquad
 \left(\frac{16p}{\gamma\eps}+1\right)
          \frac{A_t\Delta}{\gamma}\le b_*p/4.
\end{equation}
Choose an integer
\begin{equation}\label{eq:stage-number}
 M\ge\frac{32B_*K_*(T_0+3)}{c_1\Delta}.
\end{equation}
The remaining parameters are chosen sequentially in
Subsection~\ref{subsec:funding}. This order will be justified by estimates
whose constants are independent of the later stage sizes.

\subsection{Paths, replacements, and formal cells}

Start with the one-cell partition $C_0$ of the assignment cube. At stage
$j$, given the preceding partition $C_{j-1}$, use these path laws
$D^-\longrightarrow D^+$ for a sign function $\tau$ constant on its cells:
\begin{enumerate}
\item $-\tau r\longrightarrow+\tau r$ for each of the $V$ vertical laws;
\item $+\tau r\longrightarrow+\tau r'$ and
      $-\tau r'\longrightarrow-\tau r$ for each ordered pair $r\le r'$;
\item $-\tau r\longrightarrow ar$ and $ar\longrightarrow+\tau r$,
      where $r$ is iid and $a\sim\lambda$ is fresh.
\end{enumerate}
Every endpoint has absolute value at most one, and every coordinate of
$\tau(D^+-D^-)$ is nonnegative. All predicates and any amplitude in a
path are fresh samples from its specified law.

Independently in each unfrozen slot of stage $j$, retain the primary
summand with probability
$1-\alpha_j$; call this choice shared. Otherwise draw a uniform menu
entry, its fresh path data, a uniform descriptor
$\tau\in\{-1,+1\}^{J_{j-1}}$, a uniform time $u\in[0,1]$, and a fair
sign $\chi\in\{-1,+1\}$. Replace that slot by
\[
 s_j\chi[(1-u)D^-+uD^+].
\]
Here $J_{j-1}$ is the number of formal cells at the preceding level.
The score after all replacements through stage $j$ is $B^j$; future
slots and frozen slots still have their primary summands.

Construct $C_j$ by refining $C_{j-1}$ with all queried positions in the
current unfrozen primary presentations, including those in replaced
slots, and all queried positions in the current fresh paths. Use fixed
orders and pad each list to a deterministic length with coordinates
whose value is identically zero. Retain all formal binary cells,
including inconsistent and empty ones. Thus $J_j$ is a deterministic
constant once the sizes through stage $j$ are fixed. A descriptor defines
the indicated sign function on actual cells; values on empty cells do
not matter. This construction inspects neither amplitudes, times, nor
signs when forming the partition after its question lists are known.
Frozen and future query lists do not enter $C_j$.

The whole experiment first draws primary samples and freeze sets, then
performs these independent stage trials given the preceding partitions.
Trials use no lookup and no analytic reference signs. Although a
reference sign defined below can depend on future primary summands,
it never enters the generative history.

At truth, conditional on frozen variables and earlier actual trials,
the current summand differences are independent and centered. Their
second moments are at most $4s_j^2\alpha_j$. Thus
\begin{equation}\label{eq:replacement-truth}
 \E|B^j(x)-B^{j-1}(x)|
 \le\delta_j:=4\sqrt{\alpha_j/(pM)}.
\end{equation}
For each $j$, let $S_j$ be the first $m_j$ indices of $I_j$ in their
fixed order. Delete their current summands from $B^j$ to obtain
$B_\circ^j$, and delete only slot $i$ to obtain $B^j_{\neg i}$.
Conditional on the other summands and preceding history, these deleted
bases do not use the variables in the deleted slots.
The single-slot deletion will provide the background for varying that
slot's predicate. Deleting all of $S_j$ will let us choose one orientation,
from the remaining score and preceding partitions, for every tested slot.

\subsection{Smoothing and the comparison measures}

For $\xi>0$ use the probability density
\[
 f_\xi(a)=\frac{15}{16\xi}(1-a^2/\xi^2)^2\ind_{[-\xi,\xi]}(a).
\]
It is even, continuously differentiable after extension by zero, and
decreasing on $[0,\xi]$. Direct integration gives
\begin{equation}\label{eq:density-identities}
 \int|f_\xi'|=\frac{15}{8\xi},\qquad
 \int_0^\xi 2b[-f_\xi'(b)]\dd b=1.
\end{equation}
For every $z\in C_\ell$, let $\theta_{\ell,z}$ have this density with
scale $\xi_\ell$, independently. Put
$\Theta_\ell=\sum_{z\in C_\ell}\theta_{\ell,z}\ind_z$ and
$\Theta_{\le j}=\sum_{\ell=0}^j\Theta_\ell$. With partitions fixed, set
\[
 h_j(B)=\E_\theta h(B+\Theta_{\le j}),\qquad
 K_{j,z}(B)=\left.\frac{\partial}{\partial a}
                   h_j(B+a\ind_z)\right|_{a=0}.
\]
Translation of the last density proves continuous differentiability in
all cell shifts, even for merely measurable $h$. Indeed the translated
density is continuously differentiable in $L^1$, and $h$ is bounded.
Derivatives on empty cells vanish. Write $K_j=K_j(B^j)$ and
$|K_j|_1=\sum_z|K_{j,z}|$. The map $\pi$ sums fine coordinates on each
parent cell. Define the common reference descriptor
\begin{equation}\label{eq:reference-sign}
 \sigma^j=\operatorname{sign}K_{j-1}(B_\circ^j),
\end{equation}
with ties positive and using only partitions through level $j-1$.
It is common to every tested index in $S_j$.

For each mode and formal descriptor, a match in a slot means replacement
with that mode and descriptor and $\chi=+1$. Its probability is
\[
 \beta_j=\frac{\alpha_j}{2L2^{J_{j-1}}}.
\]
Mark a consistent cell in $C_j$ if any of these conditions holds:
\begin{enumerate}
\item For some mode and descriptor, the match fraction in $S_j$ exceeds
      $2\beta_j$.
\item For some vertical mode and descriptor, the fraction in $S_j$ of
      matches whose path predicate is one at the cell exceeds
      $2\beta_jp_{\rm law}$.
\item The fraction in $I_j$ that is shared and has primary predicate
      one at the cell exceeds $2p$.
\end{enumerate}
The first condition, if it occurs, marks every consistent cell. All
these rules use actual query values and trial data, and do not use $h$.
Every descriptor is checked, so on an unmarked cell the empirical bounds
also hold for the lookup-dependent descriptor $\sigma^j$. They will bound
match-weighted gradient sums without assuming that the fine gradient is
independent of the match indicators. A separate comparison form below
will charge for gradient mass on the marked cells.
At a fixed perfect witness, condition on frozen data and earlier actual
trials. Indicators in each count are independent over current slots.
Plain matches have mean $\beta_j$, vertical hit matches have mean
$\beta_jp_{\rm law}$, and shared hits have mean at most $p$.
For $m$ independent indicators of mean $b>0$, Chebyshev's inequality
bounds the probability that their average exceeds $2b$ by $1/(mb)$.
A zero mean gives a zero count almost surely. If $p_{\min}$ is the
minimum positive vertical honest probability (including $p$), a union
bound gives
\begin{equation}\label{eq:marking-truth}
 \Prob(x\text{ lies in a marked cell})
 \le\frac{2L2^{J_{j-1}}}{m_j\beta_jp_{\min}}+\frac2{n_jp}.
\end{equation}

We now specify every additional raw form. All expectations below include
the indicated score experiment; numerical noises are shared whenever
their levels appear in both arguments.

For any set $A$ of level-$j$ cells, define
\[
 N_j(B;A)=\E\sum_{z\in A}\int_0^{\xi_j}[-f_{\xi_j}'(b)]
 \big|h(B+\Theta_{\le j}^{z,+b})
           -h(B+\Theta_{\le j}^{z,-b})\big|\dd b,
\]
where the superscript fixes only coordinate $(j,z)$ and all other noises
are averaged. Translating the density and pairing $b$ with $-b$ gives
\[
 K_{j,z}(B)=\int_0^{\xi_j}[-f_{\xi_j}'(b)]
 \E_{\rm other\ noises}
 [h(B+\Theta_{\le j}^{z,+b})-h(B+\Theta_{\le j}^{z,-b})]\dd b.
\]
Thus $N_j$ dominates the expected variation on its tested cells. Its
truth cost is at most the probability that the witness cell is in $A$:
all other coordinates contribute zero, and the contributing integral
is one by Equation~\ref{eq:density-identities}. Include $N_M(B^M;C_M)$
with weight one and $N_j(B^j;\text{marked cells})$ with weight $T_0/a_j$,
where $a_j=\min(p,\Delta\beta_j)$.

For $A=B^{j-1}$ or $A=B_\circ^j$, include with weight $T_0M$ the form
\[
 P_j(A)=\E\sum_{z\in C_{j-1}}
 \int_{-\xi_{j-1}}^{\xi_{j-1}}|f_{\xi_{j-1}}'(b)|
 \big|h(B^j+\Theta_{<j}^{z,b}+\Theta_j)
              -h(A+\Theta_{<j}^{z,b})\big|\dd b.
\]
Here the fixed coordinate is $(j-1,z)$ and all other preceding noises
have their usual law. Differentiating with respect to this preceding
coordinate shifts all its children equally in the first term; its
derivative is $(\pi K_j)_z$. The second derivative is $K_{j-1,z}(A)$.
Taking absolute values under the common integral proves
\[
 \E|\pi K_j-K_{j-1}(A)|_1\le P_j(A).
\]
With $\Lambda_{j-1}=15J_{j-1}/(8\xi_{j-1})$, its unweighted truth
cost is at most $\Lambda_{j-1}(\delta_j+\xi_j)$ for the first choice
of $A$, and $\Lambda_{j-1}(m_js_j+\xi_j)$ for the second.

Include also the following ordinary coupled raw differences:
\begin{enumerate}
\item Compare $B^0$ to $B^j+\Theta_{<j}$, with weight $T_0M/e$.
\item For uniform $i\in I_j$ forced shared, compare
      $B^j+\Theta_{\le j}$ to $B^j+\Theta_{<j}$, with weight
      $T_0/(e\,p\,s_j^2)$.
\item For each mode and each of its two endpoints, force uniform
      $i\in S_j$ into that mode with independent uniform descriptor
      and $\chi=+1$. Rebuild $C_j$ using its fresh path data. For
      $B=B^j_{\neg i}+s_jD^\pm$, compare $B+\Theta_{\le j}$ to
      $B+\Theta_{<j}$, with weight $T_0/(s_j\Delta2^{-J_{j-1}})$.
\end{enumerate}
The forced trial has its prescribed conditional law
given earlier stages, and all other trials retain their ordinary laws.
The primary presentation remains in the query list. The two last types
of form have unweighted truth cost at most $\xi_j$. The first has cost
at most $\sum_{k\le j}\delta_k+\sum_{\ell<j}\xi_\ell$.

All coefficients in these measures are nonnegative. Their laws use
uniform descriptors; the lookup-dependent descriptor $\sigma^j$ in
Equation~\ref{eq:reference-sign} occurs only in the analysis.

\subsection{A sequential rational funding schedule}\label{subsec:funding}

The schedule makes all comparison truth costs simultaneously small while
retaining the variance-relative estimates needed below. Every choice uses
only parameters already fixed.

Put $\eta=1/[100M(L+1)]$ and $d_{\rm cap}=\eta e/(4T_0M^2)$.
Choose positive rational $\xi_0\le d_{\rm cap}$. At stage $j$, the
numbers $J=J_{j-1}$ and $\Lambda=\Lambda_{j-1}$ are already known.
Choose the following parameters in this order:
\begin{enumerate}
\item Choose rational $0<\alpha_j\le\min(1/2,e)$ so small that
\[
 \delta_j\le\min(d_{\rm cap},\eta/(2T_0M\Lambda)).
\]
Squaring the right-hand bound supplies a rational sufficient condition.
Set $\beta_j=\alpha_j/(2L2^J)$ and $a_j=\min(p,\Delta\beta_j)$.
\item Choose an integer
\[
 m_j\ge\frac{2T_0(2L2^J)}{\eta a_j\beta_jp_{\min}}.
\]
\item Choose an even integer
\[
 n_j\ge\max\left(2m_j,4k,\frac4{pM},\frac{4T_0}{\eta a_jp},
             \frac{16T_0^2M\Lambda^2m_j^2}{\eta^2p}\right),
\]
and rational $s_j$ satisfying Equation~\ref{eq:score-scales}. Such a
rational can be found between the two endpoints by rational interval
refinement; the bound $n_j\ge4/(pM)$ ensures $s_j\le1$.
\item Choose rational $\xi_j>0$ with
\[
 \xi_j\le\min\left(d_{\rm cap},\frac\eta{2T_0M\Lambda},
          \frac{\eta e p s_j^2}{T_0},
          \frac{\eta s_j\Delta2^{-J}}{T_0}\right).
\]
\end{enumerate}
The marked form costs at most $\eta$ by
Equation~\ref{eq:marking-truth}: its two terms cost at most $\eta/2$
each after multiplication by $T_0/a_j$. Each projected-gradient form
costs at most $\eta$; for deletion use
$m_js_j\le\eta/(2T_0M\Lambda)$, which follows from the last lower
bound on $n_j$. Each forced form costs at most $\eta$. The cumulative
demand comparison costs at most
$(T_0M/e)2Md_{\rm cap}=\eta/2$. There are $5+2L$ tunable forms per
stage, so their total truth cost is less than
$M(5+2L)\eta<1$. Together with resampling and final variation this is
less than three, and in particular at most four. All earlier truth-cost
estimates are uniform over the later primary backgrounds, so subsequent
parameter choices preserve the allocated bounds.

The consequences of Equation~\ref{eq:assumed-score-cut} are
\begin{equation}\label{eq:funded-bounds}
 \begin{gathered}
 \E|K_M|_1\le T_0\nu,\qquad
 \E\sum_{z\text{ marked}}|K_{j,z}|\le a_j\nu,\\
 \E|\pi K_j-K_{j-1}|_1\le\nu/M,\qquad
 \E|\pi K_j-K_{j-1}(B_\circ^j)|_1\le\nu/M,\\
 \E|h(B^0)-h_{j-1}(B^j)|^2\le e\nu/M.
 \end{gathered}
\end{equation}
The square bound follows by averaging raw comparisons and using
$|u-v|^2\le|u-v|$ for $u,v\in[0,1]$. In addition,
\begin{equation}\label{eq:forced-shared-bound}
 \E^{i\ {\rm uniform\ shared}}
      |h_j(B^j)-h_{j-1}(B^j)|\le e p s_j^2\nu,
\end{equation}
and each uniform-descriptor endpoint experiment satisfies
\begin{equation}\label{eq:forced-endpoint-bound}
 \E|h_j(B)-h_{j-1}(B)|\le s_j\Delta2^{-J_{j-1}}\nu.
\end{equation}

\subsection{Orientation and clipping}

The comparison measures and all numerical parameters are now fixed. It
remains to locate a stage at which a common amplitude signal survives
the accumulated approximation and clipping errors.

Put
\[
 U=\max_{0\le j\le M}\E[|K_j|_1\mid\cB],\qquad
 v=\E[U\mid\cZ].
\]
For each $j$, let $E_j=\E[|\pi K_j-K_{j-1}|_1\mid\cB]$.
The pointwise inequality $|K_{j-1}|_1\le|K_j|_1+
|\pi K_j-K_{j-1}|_1$ gives
$U\le\E[|K_M|_1\mid\cB]+\sum_jE_j$. Therefore
\begin{equation}\label{eq:U-mean}
 \E U=\E v\le(T_0+1)\nu.
\end{equation}
Define $d_{j,z}=|K_{j,z}|-\sigma^j_{\operatorname{par}(z)}K_{j,z}\ge0$.
For vectors $a,b$ and $\sigma=\operatorname{sign}b$,
$|a|_1-\sigma\cdot a\le2|a-b|_1$. Applying this with
$a=\pi K_j$, $b=K_{j-1}(B_\circ^j)$, and then the across-stage
comparison, gives
\[
 \sum_zd_{j,z}\le |K_j|_1-|K_{j-1}|_1
       +|\pi K_j-K_{j-1}|_1
       +2|\pi K_j-K_{j-1}(B_\circ^j)|_1.
\]
Summing and dropping $-|K_0|_1$ proves
\begin{equation}\label{eq:defect-telescope}
 \sum_j\E\sum_zd_{j,z}\le(T_0+3)\nu.
\end{equation}
Call $j$ good when its expected defect is at most $\Delta\nu$.

Transferring the amplitude signal requires both smoothing levels.
A current predicate $r_i$ can split a cell of $C_{j-1}$, so smoothing
along the preceding cells need not permit differentiation in its
direction. Fine smoothing $h_j$ does: $r_i$ is constant on every cell
of $C_j$. Its smoothing law, however, retains the current query partition
and can depend on the predicate's presentation. The coarse smoothing
$h_{j-1}$ omits that partition, while its score still contains the
contribution of $r_i$. Section~\ref{sec:paths} will use this omission to
express the coarse amplitude projection through a common deleted-slot
background, depending on the current presentation only through the actual
predicate. We therefore preserve the coarse component as the source of
the response and use the fine component for derivative estimates.

For $i\in I_j$, define $b_i$ as the centered amplitude component of
$h_{j-1}(B^j)$ divided by $s_j$ under the experiment forcing $i$ shared:
\[
 s_jb_i=\E^{i,\rm sh}[h_{j-1}(B^j)\mid\cB,a_i]
                 -\E^{i,\rm sh}[h_{j-1}(B^j)\mid\cB].
\]
Define $b_i^+$ by replacing $h_{j-1}$ with $h_j$, and put
$g_i=\E[b_i\mid\cZ,\cR_i,a_i]$.
Forcing shared does not alter the primary or freeze law. On replacement
the score and partition ignore the primary amplitude $a_i$. Thus the
corresponding amplitude projection under the ordinary law is exactly
$(1-\alpha_j)s_jb_i$. Conditional Bessel inequalities, applied at each
stage to the last line of Equation~\ref{eq:funded-bounds}, give
\[
 \|((1-\alpha_j)b_i-b_i^0)_i\|_\Sigma^2\le e\nu.
\]
Since $\alpha_j\le e\le1/2$ and $\|b^0\|_\Sigma\le\sqrt\nu$,
solving for $b-b^0$ gives
\begin{equation}\label{eq:component-approximations}
 \|b-b^0\|_\Sigma\le4\sqrt{e\nu},\qquad
 \|b^+-b\|_\Sigma\le\sqrt{2e\nu}.
\end{equation}
For the second bound, projection contraction in each forced experiment
bounds $s_j^2\E|b_i^+-b_i|^2$ by the expected squared difference of
the two smoothed values. Sum over $i$, use
Equation~\ref{eq:forced-shared-bound}, and use
$|I_j|s_j^2\le2/(pM)$.

Fine differentiation along the primary amplitude gives, for every $a$,
\begin{equation}\label{eq:component-gradient}
 |b_i^+(a)|\le\frac4{1-\alpha_j}
 \E\left[\ind_{\{i\ {\rm shared}\}}
             \sum_z|K_{j,z}|r_i(z)\,\middle|\,\cB\right]\le8U.
\end{equation}
To check the first inequality, bound the centered amplitude value by
the integral of absolute derivatives over the whole interval $[-1,1]$.
The derivative of $h_j$ along $a$ is $s_j\sum_zK_{j,z}r_i(z)$.
The partition is fixed along this amplitude path. Lebesgue measure on
$[-1,1]$ is bounded by $4\lambda$, giving the factor four; conditioning
a trial to be shared gives $1/(1-\alpha_j)$.
Average Equation~\ref{eq:component-gradient} over $I_j$. On unmarked
cells the shared-hit fraction is at most $2p$, and on marked cells at
most one. Equations~\ref{eq:funded-bounds} and~\ref{eq:U-mean} imply
\begin{equation}\label{eq:component-L1}
 \E\operatorname{avg}_{i\in I_j}|b_i^+|
 \le16p(T_0+1)\nu+8p\nu\le K_*p\nu.
\end{equation}

We now clip the fine component to obtain a bounded signal. Set
$A=\{U\le B_*,\ v\le B_*\}$, which is $\cB$-measurable, and
define $\widetilde b_i=b_i^+\ind_A$ and
$\widetilde g_i=\E[\widetilde b_i\mid\cZ,\cR_i,a_i]$.

\begin{table}[htbp]
\centering
\small
\begin{tabular}{@{}lp{0.72\textwidth}@{}}
\toprule
Notation & Definition and role \\
\midrule
$b_i^0,\ g_i^0$ & Centered amplitude component of the demand lookup
$h(B^0)$, divided by $s_j$, and its projection onto
$(\cZ,\cR_i,a_i)$; these supply the initial signal. \\[3pt]
$b_i,\ g_i$ & Component of the coarse lookup $h_{j-1}(B^j)$ when
$i$ is forced shared, and the same conditional projection;
Section~\ref{sec:paths} expresses $g_i$ through the common deleted-slot
background. \\[3pt]
$b_i^+$ & Component with fine smoothing $h_j$ in the same forced
experiment; differentiation gives the pointwise and first-moment bounds. \\[3pt]
$\widetilde b_i,\ \widetilde g_i$ & The clipped component
$b_i^+\ind_{\{U\le B_*,\,v\le B_*\}}$ and its conditional projection;
these give bounded quantities for the following norm comparisons. \\
\bottomrule
\end{tabular}
\caption{Amplitude components. The fine and clipped quantities retain a
signal in the coarse projection $g_i$.}
\label{tab:amplitude-notation}
\end{table}

The conditional cap in Equation~\ref{eq:amplitude-signal}, even with
its weaker upper bound one, and Markov's inequality give
\[
 \|b^0\ind_{A^c}\|_\Sigma^2
 \le\Prob(A^c)\le2(T_0+1)\nu/B_*.
\]
Projection contraction, Equation~\ref{eq:component-approximations},
and the triangle inequality now give
\begin{align}
 \|\widetilde g-g^0\|_\Sigma
 &\le[6\sqrt e+\sqrt{2(T_0+1)/B_*}]\sqrt\nu,
       \label{eq:clipped-signal-distance}\\
 \|\widetilde g-g\|_\Sigma
 &\le[10\sqrt e+\sqrt{2(T_0+1)/B_*}]\sqrt\nu.
       \label{eq:clipped-comparison-distance}
\end{align}
For every family $q_i$, Equation~\ref{eq:score-scales} implies
\begin{equation}\label{eq:norm-conversion}
 \frac p2\|q\|_\Sigma^2
 \le\operatorname{avg}_j\E\operatorname{avg}_{i\in I_j}|q_i|^2
 \le2p\|q\|_\Sigma^2.
\end{equation}
The first bracket above is at most $\sqrt c/25$ and the second at
most $\sqrt c/20$ by Equation~\ref{eq:preconjunction-constants}.
Using the amplitude signal and Equation~\ref{eq:norm-conversion},
we obtain
\begin{equation}\label{eq:clipped-signal}
 \operatorname{avg}_j\E\operatorname{avg}_{i\in I_j}
       |\widetilde g_i|^2\ge c_1p\nu,\qquad
 \operatorname{avg}_j\E\operatorname{avg}_{i\in I_j}
       |\widetilde g_i-g_i|^2\le\tfrac1{10}c_1p\nu.
\end{equation}
Also $|\widetilde g_i|\le8B_*$, its averaged first moment is at most
$K_*p\nu$ at each stage by Equation~\ref{eq:component-L1}, and it
vanishes when $v>B_*$.

At most $(T_0+3)/\Delta$ stages are bad. Each contributes at most
$8B_*K_*p\nu$ to the unaveraged second moment of $\widetilde g$.
Equation~\ref{eq:stage-number} bounds their average contribution by
$c_1p\nu/4$. On $|g_i|<\gamma$, the elementary inequality
\[
 |\widetilde g_i|^2\le2\gamma|\widetilde g_i|
                              +|\widetilde g_i-g_i|^2
\]
holds, as follows by expanding the last square and using
$|g_i|<\gamma$. Its average is at most
$(1/8+1/10)c_1p\nu$. Thus at least $c_1p\nu/2$ of the average
second moment remains on good stages with $|g_i|\ge\gamma$ and
$v\le B_*$. The bound $|\widetilde g_i|^2\le64B_*^2$ proves that
some good stage $j$ satisfies
\begin{equation}\label{eq:good-stage-signal}
 \Prob_{\cZ,\,i\ {\rm uniform\ in}\ I_j,\,\cR_i,a_i}
             (|g_i|\ge\gamma,\ v\le B_*)\ge b_*p\nu.
\end{equation}
We fix this stage for the rest of the analytic argument.
